\documentclass[12pt]{article}

\usepackage[utf8]{inputenc}
\usepackage[T1]{fontenc}
\usepackage{amsmath, amssymb}
\usepackage{graphicx}
\usepackage{float}
\usepackage{listings}
\usepackage{xcolor}
\usepackage[top=1in, bottom=1in, left=1in, right=1in]{geometry}
\usepackage{titling}
\setlength{\droptitle}{-3em}

\lstset{
    language=Python,
    basicstyle=\ttfamily\small,
    keywordstyle=\color{blue},
    commentstyle=\color{green!60!black},
    stringstyle=\color{red},
    numbers=left,
    numberstyle=\tiny\color{gray},
    frame=single,
    backgroundcolor=\color{gray!10},
    breaklines=true,
    showstringspaces=false
}

\lstdefinestyle{plainschema}{
    language=,
    basicstyle=\ttfamily\small,
    frame=single,
    backgroundcolor=\color{gray!10},
    breaklines=true,
    showstringspaces=false,
    numbers=none
}

\lstdefinestyle{sqlstyle}{
    language=SQL,
    basicstyle=\ttfamily\small,
    keywordstyle=\color{blue},
    commentstyle=\color{green!60!black},
    stringstyle=\color{red},
    frame=single,
    backgroundcolor=\color{gray!10},
    breaklines=true,
    showstringspaces=false,
    numbers=none
}

\geometry{margin=1in}

\title{Database Notes: Schemas, Keys, and Schema Design}
\author{
Greyson Knippel \\[4pt]
CPSC 321 --- Gonzaga University
}
\date{\today}

\begin{document}

\maketitle

\noindent Reference notes for an intro database course. Covers relational schema
vocabulary, the key hierarchy, how to derive keys from stated constraints, and
foreign key placement in schema design.

\section{Relation Schema}

A \textbf{relation schema} has three parts:

\begin{enumerate}
    \item \textbf{Relation name} --- e.g. \texttt{STUDENT}
    \item \textbf{Attributes} --- $A_1, A_2, \ldots, A_n$ (the columns)
    \item \textbf{Domains} --- $\mathrm{dom}(A_i)$, the allowed atomic values for each attribute
\end{enumerate}

Notation: $R(A_1, A_2, \ldots, A_n)$

\subsection*{Example}

\begin{lstlisting}[style=plainschema]
STUDENT(Name, SSN, Age, Major)
\end{lstlisting}

\begin{table}[H]
\centering
\begin{tabular}{|l|l|c|l|}
\hline
\textbf{Name} & \textbf{SSN} & \textbf{Age} & \textbf{Major} \\
\hline
Ana Ruiz & 305-61-2435 & 19 & CS \\
Ben Cho & 422-11-9087 & 21 & Math \\
Dee Patel & 700-53-1122 & 20 & CS \\
\hline
\end{tabular}
\end{table}

\begin{table}[H]
\centering
\begin{tabular}{|l|l|l|}
\hline
\textbf{Term} & \textbf{Meaning} & \textbf{Here} \\
\hline
Relation name & the table's name & \texttt{STUDENT} \\
Attribute & a column & \texttt{Name}, \texttt{SSN}, \texttt{Age}, \texttt{Major} \\
Domain & allowed values for an attribute & $\mathrm{dom}(\text{Age}) = $ int 1--120 \\
Degree / arity & number of attributes & 4 \\
Tuple & one row & \texttt{<Ana Ruiz, 305-61-2435, 19, CS>} \\
Cardinality & number of tuples & 3 \\
Attribute value & one cell, must be \textbf{atomic} & $t[\text{Age}] = 19$ \\
Relation state $r(R)$ & the whole set of tuples right now & these 3 rows \\
\hline
\end{tabular}
\end{table}

\textbf{Schema vs.\ instance:} the schema is the definition (intension); the
relation state/instance is the data at a moment (extension).

\textbf{Relation schema vs.\ relational database schema:} the latter is
\emph{all} relation schemas plus integrity constraints.

A relation is a \textbf{set of tuples} --- no duplicates, no row order. Values
\emph{within} a tuple are ordered, matched positionally to the attributes.

\section{The Key Hierarchy}

\subsection*{Definitions}

\begin{itemize}
    \item \textbf{Superkey} --- any attribute set that uniquely identifies a tuple. Extras allowed.
    \item \textbf{Candidate key} --- a superkey that is \textbf{minimal}: removing \emph{any} attribute breaks uniqueness.
    \item \textbf{Primary key} --- the candidate key you choose as the official identifier. Never NULL.
    \item \textbf{Alternate key} --- the candidate keys you didn't choose.
    \item \textbf{Foreign key} --- attribute(s) referencing another relation's primary key.
    \item \textbf{Surrogate key} --- a meaningless system-generated ID (auto\_increment, UUID).
\end{itemize}

\subsection*{Simple vs.\ composite --- just a count}

\begin{itemize}
    \item \textbf{Simple key} = one attribute
    \item \textbf{Composite key} = two or more attributes taken together
\end{itemize}

Orthogonal to primary/candidate. A composite key can absolutely be the primary key.

\subsection*{The Minimality Test}

\begin{quote}
Remove each attribute \textbf{one at a time}. If \textbf{EVERY} removal breaks
uniqueness, it's a candidate key. If \textbf{ANY} removal survives, it's a
superkey only.
\end{quote}

\textbf{Breaking is the PASS condition.} If it still works after removal, that
attribute was dead weight and the set was not minimal.

One-line version: \textbf{a candidate key is a superkey you cannot make any smaller.}

\subsection*{Worked Example --- EMPLOYEE}

\begin{lstlisting}[style=plainschema]
EMPLOYEE(EmpID, Email, Phone, Name, Dept)
\end{lstlisting}

Constraints: \texttt{EmpID}, \texttt{Email}, \texttt{Phone} each unique.
\texttt{Name}, \texttt{Dept} repeat.

\begin{table}[H]
\centering
\begin{tabular}{|l|c|c|l|}
\hline
\textbf{Set} & \textbf{Superkey?} & \textbf{Candidate key?} & \textbf{Why} \\
\hline
\texttt{\{EmpID\}} & \checkmark & \checkmark & unique + minimal \\
\texttt{\{Email\}} & \checkmark & \checkmark & unique + minimal \\
\texttt{\{EmpID, Name\}} & \checkmark & $\times$ & remove \texttt{Name} $\to$ still works $\to$ not minimal \\
\texttt{\{EmpID, Email\}} & \checkmark & $\times$ & two keys stapled together --- remove either, still works \\
\texttt{\{Name, Dept\}} & $\times$ & $\times$ & not unique \\
all 5 attributes & \checkmark & $\times$ & the whole row is always a superkey \\
\hline
\end{tabular}
\end{table}

Counts: \textbf{28 superkeys} ($2^5 = 32$ subsets, minus the 4 subsets of
\texttt{\{Name, Dept\}}), but only \textbf{3 candidate keys}.

\begin{quote}
Gotcha: \texttt{\{EmpID, Email\}} is the \emph{worst} offender, not the best ---
it contains two complete keys when one would do. Keyring analogy: two keys that
open the same door. A candidate key is a keyring where losing \textbf{any} key
locks you out.
\end{quote}

\subsection*{Worked Example --- CAR (composite candidate key)}

\begin{lstlisting}[style=plainschema]
CAR(LicensePlate, State, VIN, Make, Model, Year)
\end{lstlisting}

License plates are unique only \emph{within} a state.

\begin{itemize}
    \item Candidate keys: \texttt{\{VIN\}} (simple) and \texttt{\{LicensePlate, State\}} (composite)
    \item \texttt{\{LicensePlate\}} alone is \textbf{not} a superkey --- the trap
\end{itemize}

A relation can have a simple and a composite candidate key at the same time.

\subsection*{Worked Example --- Derive From Data}

\begin{lstlisting}[style=plainschema]
R(A, B, C)
\end{lstlisting}

\begin{table}[H]
\centering
\begin{tabular}{|c|c|c|}
\hline
\textbf{A} & \textbf{B} & \textbf{C} \\
\hline
a1 & b1 & c1 \\
a1 & b2 & c1 \\
a2 & b1 & c2 \\
a2 & b2 & c2 \\
\hline
\end{tabular}
\end{table}

\begin{itemize}
    \item \texttt{\{A\}}, \texttt{\{B\}}, \texttt{\{C\}} --- all have duplicates $\times$
    \item \texttt{\{A,B\}} --- distinct $\to$ superkey, and minimal $\to$ \textbf{candidate key} \checkmark
    \item \texttt{\{A,C\}} --- (a1,c1) twice $\times$ (because $A \to C$, so $C$ adds nothing)
    \item \texttt{\{B,C\}} --- distinct $\to$ \textbf{candidate key} \checkmark
    \item \texttt{\{A,B,C\}} --- superkey, not minimal
\end{itemize}

\textbf{3 superkeys, 2 candidate keys.}

\begin{quote}
\textbf{Warning:} An instance can only DISPROVE a key, never PROVE one. Keys are
constraints over all possible states. Uniqueness in one snapshot may be
coincidence. Keys come from the rules of the domain, not from staring at data.
\end{quote}

\subsection*{The Minimality Trap}

\begin{lstlisting}[style=plainschema]
ENROLLMENT(StudentID, CourseID, Semester, Grade)
\end{lstlisting}

All four columns together are unique --- but drop \texttt{Grade} and it still
works, so it's a \textbf{superkey only}. Candidate key is
\texttt{\{StudentID, CourseID, Semester\}}.

Rule of thumb: \textbf{if you can delete an attribute from your answer and it
still works, your answer is wrong.}

\section{Deriving Keys From Stated Constraints}

Word problems state constraints in the form:

\begin{quote}
``at most one \textbf{[B]} for a given \textbf{[A]}''
\end{quote}

This gives you a \textbf{functional dependency}: $A \to B$.

\subsection*{The Rule}

\begin{enumerate}
    \item Translate the sentence into an FD.
    \item Ask: \textbf{does the left side determine EVERY remaining attribute?}
    \begin{itemize}
        \item Yes $\to$ the left side is the key.
        \item No $\to$ add the uncovered attributes, then re-run the minimality test.
    \end{itemize}
\end{enumerate}

\subsection*{Three Worked Cases}

\textbf{Case 1 --- \texttt{ALLOWED\_PLAN(vt\_id, p\_id)}}

``A vehicle type has at most one plan; a plan serves many vehicle types.''
\begin{itemize}
    \item FD: $vt\_id \to p\_id$
    \item Left side covers everything else \checkmark
    \item \textbf{PK = \{vt\_id\}} (simple)
    \item Both columns are also FKs. Since it's 1:N, this table could be folded
    into \texttt{VEHICLE\_TYPE} as an FK column --- junction tables earn their
    keep only for M:N.
\end{itemize}

\textbf{Case 2 --- \texttt{SPY(agent, spies\_on, spies\_for)}}

``A country has at most one agent that spies on a given country.''
\begin{itemize}
    \item FD: $\{spies\_on, spies\_for\} \to agent$
    \item Left side covers everything else \checkmark
    \item \textbf{PK = \{spies\_on, spies\_for\}} (composite)
    \item \texttt{agent} is not in the key --- a row is an \emph{assignment}, not
    a person. One agent can hold many assignments (multiple targets, multiple
    employers, double agent).
\end{itemize}

\textbf{Case 3 --- \texttt{SPY(agent, spies\_on, spies\_for)}, different constraint}

``An agent spies on at most one country.''
\begin{itemize}
    \item FD: $agent \to spies\_on$
    \item Left side does \textbf{not} cover \texttt{spies\_for} $\times$ $\to$ patch the hole
    \item \textbf{PK = \{agent, spies\_for\}}
    \item Note: \texttt{agent} is only \emph{part} of the key but determines
    \texttt{spies\_on} $\to$ \textbf{partial dependency $\to$ 2NF violation}.
    Decomposes into \texttt{AGENT(agent, spies\_on)} and
    \texttt{ASSIGNMENT(agent, spies\_for)}.
\end{itemize}

\textbf{Same three columns, same table shape --- the constraint decides the key.}

\section{Foreign Keys}

\begin{lstlisting}[style=plainschema]
EMPLOYEE(EmpID, Name, DeptID)
                        --FK--> DEPARTMENT(DeptID)
\end{lstlisting}

\textbf{Referential integrity:} every FK value must either match an existing PK
value in the parent table, or be NULL.

\subsection*{Three Things Students Get Wrong}

\begin{enumerate}
    \item \textbf{An FK does NOT have to be unique.} Many employees, one
    department --- \texttt{DeptID} repeats. Uniqueness is required on the
    \emph{parent} side only.
    \item \textbf{An FK CAN be NULL} (unless separately declared \texttt{NOT NULL}). A PK never can.
    \item \textbf{Direction matters.} The FK lives on the \textbf{many} side.
\end{enumerate}

\subsection*{What It Prevents}

\begin{lstlisting}[style=sqlstyle]
INSERT INTO EMPLOYEE VALUES ('E5','Kim','D9');  -- REJECTED: D9 doesn't exist (orphan)
DELETE FROM DEPARTMENT WHERE DeptID = 'D1';     -- REJECTED: employees still reference D1
\end{lstlisting}

\subsection*{Referential Actions}

\begin{table}[H]
\centering
\begin{tabular}{|l|l|}
\hline
\textbf{Action} & \textbf{On delete of the parent row} \\
\hline
\texttt{RESTRICT} / \texttt{NO ACTION} & refuse the delete (default) \\
\texttt{CASCADE} & delete the children too \\
\texttt{SET NULL} & keep children, null out their FK \\
\texttt{SET DEFAULT} & set FK to a declared default \\
\hline
\end{tabular}
\end{table}

\begin{lstlisting}[style=sqlstyle]
FOREIGN KEY (DeptID) REFERENCES DEPARTMENT(DeptID)
  ON DELETE SET NULL
  ON UPDATE CASCADE
\end{lstlisting}

\subsection*{Variations}

\begin{itemize}
    \item \textbf{Composite FK} --- if the parent PK is composite, the FK must be too, all columns together.
    \item \textbf{Self-referencing FK} --- \texttt{EMPLOYEE(EmpID, Name, ManagerID)}
    where \texttt{ManagerID} $\to$ \texttt{EMPLOYEE(EmpID)}. The top of the chain
    has \texttt{ManagerID = NULL}, which is exactly why FKs must allow NULL.
    \item \textbf{PK that is also an FK} --- normal in junction tables and 1:1 relationships.
\end{itemize}

\section{Schema Design --- Where the FK Goes}

\begin{table}[H]
\centering
\begin{tabular}{|l|l|}
\hline
\textbf{Relationship} & \textbf{Where the FK goes} \\
\hline
\textbf{1:N} & FK column on the \textbf{many} side --- no extra table \\
\textbf{M:N} & \textbf{New junction table}, PK = both FKs combined \\
\textbf{1:1} & FK on either side, plus a UNIQUE constraint \\
\hline
\end{tabular}
\end{table}

\textbf{Cardinality decides this, not table shape.} Two two-column tables can
look identical and need opposite designs.

\subsection*{M:N Example --- Books and Authors}

``A book can have multiple authors, and an author can write multiple books.''

\begin{lstlisting}[style=plainschema]
BOOK(isbn, title, pub_year)
     PK: {isbn}

AUTHOR(author_id, first_name, last_name)
       PK: {author_id}

WROTE(isbn, author_id)
      PK: {isbn, author_id}
      FK: isbn      -> BOOK(isbn)
      FK: author_id -> AUTHOR(author_id)
\end{lstlisting}

\textbf{BOOK and AUTHOR have no foreign keys.} Both FKs live in the junction table.

Why the third table is mandatory: putting \texttt{author\_id = "A1, A2, A3"} in
BOOK violates \textbf{atomicity (1NF)}; \texttt{author1/author2/author3} columns
cap the author count and fill with NULLs. The junction table makes each
authorship its own row.

\begin{lstlisting}[style=sqlstyle]
CREATE TABLE WROTE (
  isbn       CHAR(13),
  author_id  INTEGER,
  PRIMARY KEY (isbn, author_id),
  FOREIGN KEY (isbn)      REFERENCES BOOK(isbn)        ON DELETE CASCADE,
  FOREIGN KEY (author_id) REFERENCES AUTHOR(author_id) ON DELETE CASCADE
);
\end{lstlisting}

Relationship attributes (e.g. \texttt{author\_position} for cover order) belong
in the junction table --- they describe the pairing, not either entity, and
don't join the key.

\subsection*{1:N Example --- Pets and Owners}

``Each pet belongs to exactly one owner.''

\begin{lstlisting}[style=plainschema]
OWNER(owner_id, name, phone)
      PK: {owner_id}

PET(pet_id, name, species, owner_id)
    PK: {pet_id}
    FK: owner_id -> OWNER(owner_id)   [NOT NULL]
\end{lstlisting}

Two tables, no junction table. A \texttt{HAS\_PET} junction table would
\emph{permit} a pet with two owners --- \textbf{your schema should make illegal
states impossible.}

Don't put \texttt{pet\_id} in OWNER --- that caps each owner at one pet.

\subsection*{The ``Exactly'' vs.\ ``At Most'' Distinction}

\begin{quote}
\textbf{``exactly one''} $\to$ FK is \textbf{NOT NULL} \\
\textbf{``at most one''} $\to$ FK \textbf{allows NULL}
\end{quote}

Graders look for this. It's also the one real argument for keeping a separate
table in a 1:N case: if most rows would have no match, a separate table avoids
a column full of NULLs.

\section{OLTP vs.\ OLAP}

\begin{table}[H]
\centering
\begin{tabular}{|l|l|l|}
\hline
 & \textbf{OLTP} & \textbf{OLAP} \\
\hline
Purpose & runs the business & analyzes the business \\
Workload & many small short transactions & few huge complex read queries \\
Query touches & a few rows, by key & millions of rows, scanned + aggregated \\
Read/write & write-heavy, high concurrency & read-mostly, bulk-loaded \\
Data & current, detailed, updated in place & historical, summarized, non-volatile \\
Design & \textbf{normalized} (3NF) --- cheap writes & \textbf{denormalized} --- cheap joins \\
Optimizes for & ACID, latency, throughput & scan performance \\
\hline
\end{tabular}
\end{table}

Why they're separate databases: a 40-minute analytical scan shouldn't lock up
order processing, and the two workloads want physically opposite designs.

\subsection*{Warehouse Vocabulary}

\textbf{Data warehouse} (Inmon's definition --- often tested verbatim):
\emph{subject-oriented, integrated, time-variant, non-volatile.}

\begin{itemize}
    \item \emph{Subject-oriented} --- organized around subjects, not applications
    \item \emph{Integrated} --- merged from many sources, formats reconciled
    \item \emph{Time-variant} --- every record carries a time dimension; history kept
    \item \emph{Non-volatile} --- loaded and read, not updated or deleted
\end{itemize}

\begin{itemize}
    \item \textbf{Data mart} --- a warehouse scoped to one department
    \item \textbf{ETL} --- Extract, Transform, Load (ELT = load raw, transform in the warehouse)
    \item \textbf{Star schema} --- central \textbf{fact table} (numeric measures + FKs)
    surrounded by denormalized \textbf{dimension tables} (Product, Store, Date, Customer)
    \item \textbf{Snowflake schema} --- star schema with normalized dimensions
    \item \textbf{Fact vs.\ dimension} --- facts are measures; dimensions are how you slice them
\end{itemize}

\subsection*{OLAP Cube Operations}

\begin{table}[H]
\centering
\begin{tabular}{|l|l|}
\hline
\textbf{Operation} & \textbf{Meaning} \\
\hline
Roll-up & aggregate up a hierarchy (day $\to$ month $\to$ year) \\
Drill-down & the reverse, go finer \\
Slice & fix one dimension to a single value \\
Dice & subcube from ranges on several dimensions \\
Pivot & rotate the axes of the view \\
\hline
\end{tabular}
\end{table}

\section*{Quick Reference Card}

\textbf{Superkey} = unique. \\
\textbf{Candidate key} = unique AND nothing removable. \\
\textbf{Primary key} = the candidate key you picked. \\
\textbf{Composite} = 2+ attributes. \textbf{Simple} = 1 attribute.

\textbf{Minimality test:} remove each attribute one at a time --- \emph{every}
removal must break it.

\textbf{Constraint $\to$ key:} turn the sentence into an FD, check whether the
left side covers all remaining attributes, patch if not.

\textbf{FK placement:} 1:N $\to$ many side. M:N $\to$ junction table with
composite PK.

\textbf{``Exactly one''} $\to$ NOT NULL. \textbf{``At most one''} $\to$ NULL allowed.

\textbf{Instances disprove keys; they never prove them.}

\end{document}
